\documentclass{beamer}
\usepackage[numberedbib]{apacite}
\usepackage[english]{babel}
\usepackage{graphicx}
\usepackage{multimedia}
\usepackage{hyperref}

\newcommand{\putat}[3]{\begin{picture}(0,0)(0,0)\put(#1,#2){#3}\end{picture}}
% Choose how your presentation looks.
%
% For more themes, color themes and font themes, see:
% http://deic.uab.es/~iblanes/beamer_gallery/index_by_theme.html
%
\mode<presentation>
{
  \usetheme{default}      % or try Darmstadt, Madrid, Warsaw, ...
  \usecolortheme{default} % or try albatross, beaver, crane, ...
  \usefonttheme{default}  % or try serif, structurebold, ...
  \setbeamertemplate{navigation symbols}{}
  \setbeamertemplate{caption}[numbered]
} 

\usepackage[english]{babel}
\usepackage[utf8x]{inputenc}

\title[Simultaneous Evolution of Morphology and Locomotion of Soft Robots by Novelty Search]{Simultaneous Evolution of Morphology and Locomotion of Soft Robots by Novelty Search}
\author{Georgios Methenitis}
\institute{University of Amsterdam}
\date{\today}

\begin{document}

\begin{frame}
  \titlepage
\end{frame}

% Uncomment these lines for an automatically generated outline.
%\begin{frame}{Outline}
%  \tableofcontents
%\end{frame}

\section{Introduction}
\begin{frame}{Introduction}
\begin{block}{Soft Robots}
\begin{itemize}
\item Inspired by nature
\item Completely soft bodies
\item Capable of developing new kinds of locomotion
\end{itemize}
\end{block}
\includegraphics[width=0.3\textwidth,height=0.25\textheight]{figures/soft_robotics_figure.png}\	
\includegraphics[width=0.3\textwidth,height=0.25\textheight]{figures/ExplodingRobot.jpg}\	
\includegraphics[width=0.3\textwidth,height=0.25\textheight]{figures/hillerPressureChamber.png}\\
\vspace{0.3cm}
Soft robots can be actuated through air pressure tubes, environmental changes ( temperature, pressure ), even explosions.
\end{frame}

\begin{frame}{Related Material}
\begin{block}{\includegraphics[scale=0.35]{figures/voxcad_logo.png}\	VoxCad Simulator~\cite{hiller2012dynamic}}
\begin{itemize}
\item Created by Jonathan Hiller and Hod Lipson
\item Voxel modeling and analyzing software
\item Capable of developing new kinds of locomotion
\end{itemize}
\end{block}
\putat{95}{-115}{\includegraphics[scale=0.12]{figures/voxcad.png}}
\begin{itemize}
\item Lattice
\item Voxels
\item Structure
\item Materials
\end{itemize}
\end{frame}

\begin{frame}[allowframebreaks]{Related Work}
\textit{Evolving virtual creature}s~\cite{sims1994evolving}
\begin{itemize}
\item Rigid body parts, joints
\item Evolution of the morphology and the control
\end{itemize}
\includegraphics[scale=0.5]{figures/evolvingVirtualCreatures.png}\\
\textit{Evolving a diversity of virtual creatures through novelty search and local competition} ~\cite{lehman2011evolving}
\begin{itemize}
\item Same experiment
\item Novelty $<$ Fitness
\item Novelty search with global competition has the best average fitness.
\end{itemize}
\newpage
\textit{Evolving soft robots with multiple materials and a powerful generative encoding.}~\cite{cheney2013unshackling}
\begin{itemize}
\item Generative encoding, Compositional pattern-producing network, CPPN.
\item Neuroevolution of augmenting topologies, NEAT.
\end{itemize}
\vspace{0.3cm}
\includegraphics[width=0.3\textwidth,height=0.25\textheight]{figures/unshacklingEvolutionFigure1.png}\	\	\	
\includegraphics[width=0.3\textwidth,height=0.25\textheight]{figures/unshacklingEvolutionFigure2.png}\	\	\	
\includegraphics[width=0.3\textwidth,height=0.25\textheight]{figures/unshacklingEvolutionFigure3.png}
\end{frame}



\begin{frame}{Compositional pattern-producing network~\cite{stanley2007compositional}}
\begin{itemize}
\item Similar to artificial neural networks
\item Different set of activation functions
\begin{figure}
\begin{center}
\includegraphics[scale=0.25]{figures/cppnNetwork.png}
\end{center}
\end{figure}
\item Produce symmetrical and repetitive patterns
\item Appropriate for problems with geometrical structure (Board games).
\end{itemize}
\end{frame}


\begin{frame}{NeuroEvolution through Augmented Topologies (NEAT)~\cite{stanley2002evolving}}
\begin{block}{Some key points of this method are:}
\begin{itemize}
\item Evolving neural network topologies along with weights
\item Crossover between different topologies
\item Structural innovation through speciation (New species have time to improve)
\end{itemize}
\end{block}
\end{frame}

\begin{frame}{Research Topics}
\begin{itemize}
\item Gravity
\begin{itemize}
\item Performance under different conditions of gravity
\end{itemize}
\item Novelty search
\begin{itemize}
\item Performance, in respect to the original fitness
\item Performance, in behavior space
\item Behavior, what is a good behavior metric?
\end{itemize}
\item Other evolutionary algorithms
\begin{itemize}
\item Genetic algorithm with direct coding
\item Random generative encoding
\item CMA-ES
\item jDE
\end{itemize}
\item Can we evolve CPPNs with other evolutionary algorithms?
\end{itemize}
\end{frame}

\begin{frame}{Things completed so far\ldots}
\begin{itemize}
%\item All tools documented
\item Replication of the results from~\cite{cheney2013unshackling}
\item Generative random encoding
\item Simple genetic algorithm
\item Own implementation of CPPN-NEAT experiment (HyperNEAT~C++~library)
\item Novelty search
\item Competition between species (novelty, fitness)
\end{itemize}
\end{frame}

\begin{frame}{Generative Random Encoding}
Only two parameters can change in this encoding.
\begin{enumerate}
\item The probability of adding a new voxel into the structure.
\item The probability that the new voxel introduced will use the same kind of material as its connection.
\end{enumerate}
We start with a random voxel inserted into the lattice.
\begin{itemize}
\item Choose whether a new voxel is going to be inserted.
\item Choose randomly its connection.
\item Choose its material.
\item Iterate
\end{itemize}
\end{frame}

\begin{frame}{Simple genetic algorithm}
\begin{itemize}
\item Each genome is represented by a stream of real numbers in $[ 0,1 ]$.
\item The length of this stream is equal to:
\begin{equation*}
l = n \times (m + 1)
\end{equation*}
, where $n$, is the number of total voxels and $m$ is the number of materials.
\item For a lattice's dimensions of $10 \times 10 \times 10$, the length is $5000$.
\item Simple genetic algorithm fails to produce locomotion.
\item No structure knowledge.
\end{itemize}
\end{frame}

\begin{frame}[allowframebreaks]{References}
\bibliographystyle{apacite}
\bibliography{ref}
\end{frame}


\end{document}
